\SourceLeaf{416}%
sich ergebenden Werthe alle dieselben Zeichen haben, wie respective die Sinus von \(\vartheta,\vartheta',\vartheta''\).

Mit Stillschweigen darf ich auch nicht übergehen, dass die gegebene allgemeine Auflösung der Aufgabe in singulären Fällen entweder ihre Anwendbarkeit ganz verliert, oder doch einiger Modification bedarf, beschränke mich hier aber nur auf eine Andeutung der erheblichsten Punkte:

I. Ist einer der beobachteten Winkel gleich \(=0\) oder \(=180^0\), so ist vortheilhaft, den betreffenden Beobachtungsplatz, auch wenn er zwischen den beiden andern liegen sollte, als \((1)\) oder \((3)\) anzunehmen. Wählet man das Letztere, so bleiben alle Theile der allgemeinen Auflösung gültig, indem man nur \(b\) als unendlich gross betrachtet, und als Zeichen in den Rechnungen beibehält, welches hernach in allen Resultaten aber von selbst wegfällt. An mehr als Einem Punkte darf aber offenbar der Winkel nicht \(0\) oder \(180^0\) sein, weil dies nur stattfindet, wenn alle drei in der Linie \((4),(5)\) liegen, wo die Aufgabe unbestimmt wäre.

II. Sind unter den beobachteten Winkeln zwei gleiche, so fallen von den Punkten \((1^*),(2^*),(3^*)\) zwei zusammen, oder eine der Grössen \(a,b,a+b\) wird \(=0\), daher auch \(fF=0\); in diesem Falle geht demnach die cubische Gleichung in eine quadratische über. Übrigens sieht man leicht, dass das Verschwinden des ersten Coëfficienten der cubischen Gleichung \textit{nur} in dem Falle der Gleichheit zweier Winkel eintritt.

III. Die gegebene Auflösung ist nicht anzuwenden, wenn die Grössen \(A,B\) den \(a,b\) proportional sind, also \(\lambda=0\) wird. In diesem Falle ist die cubische Gleichung mit Unrecht herangezogen und enthält eine der Sache fremde Wurzel, die richtige aber zweimal. Es ist nemlich klar, dass dann die beiden Combinationen, durch welche aus \([2]\) und \([3]\) die Gleichungen \([4]\) und \([5]\) abgeleitet wurden, \textit{nicht verschieden} sind; diese Gleichungen werden daher identisch, und jede für sich gibt \(x=\frac A{2a}=\frac B{2b}\). Offenbar muss dann aber eine der beiden Gleichungen \([2],[3]\) noch ferner gebraucht werden, aus deren Combination mit \([1]\) sich leicht ergibt: \(cx=(t-n)\BareRoot{(1-xx)}\). Es erhellt daraus, dass die Linie \(\mathrm{(I)}\) entweder durch den Punkt \((4)\) oder durch \((5)\) geht, und es gibt
